\documentclass[12pt, a4paper]{article}

% Pacchetti
\usepackage[utf8]{inputenc}
\usepackage[T1]{fontenc}
\usepackage[italian]{babel}
\usepackage{geometry}
\usepackage{titlesec}
\usepackage{enumitem}
\usepackage{fancyhdr}
\usepackage{xcolor}
\usepackage{booktabs}
\usepackage{array}
\usepackage{parskip}
\usepackage{hyperref}

% Margini
\geometry{top=2.5cm, bottom=2.5cm, left=3cm, right=3cm}

% Colori
\definecolor{primary}{RGB}{30, 60, 120}
\definecolor{lightgray}{RGB}{245, 245, 245}
\definecolor{darkgray}{RGB}{80, 80, 80}

% Stile titoli sezioni
\titleformat{\section}
  {\large\bfseries\color{primary}}
  {\thesection.}{0.5em}{}[\titlerule]

\titleformat{\subsection}
  {\normalsize\bfseries\color{darkgray}}
  {\thesubsection.}{0.5em}{}

% Header e Footer
\pagestyle{fancy}
\fancyhf{}
\rhead{\small\textcolor{darkgray}{Verbale di Riunione}}
\lhead{\small\textcolor{darkgray}{MirrorFish Trading}}
\rfoot{\small\textcolor{darkgray}{\thepage}}
\renewcommand{\headrulewidth}{0.4pt}

% Hyperlink
\hypersetup{
  colorlinks=true,
  linkcolor=primary,
  urlcolor=primary
}

% ─────────────────────────────────────────────
\begin{document}

% Intestazione / Titolo
\begin{center}
  {\LARGE\bfseries\color{primary} VERBALE DI RIUNIONE}\\[0.4cm]
  {\large Progetto: \textit{MirrorFish Trading}}\\[0.2cm]
  \textcolor{darkgray}{\rule{0.6\textwidth}{0.4pt}}
\end{center}

\vspace{0.5cm}

% Tabella riepilogativa
\begin{center}
\begin{tabular}{>{\bfseries}p{4cm} p{9cm}}
\toprule
Data & 14 aprile 2026 \\
\midrule
Luogo & Discord (chiamata vocale) \\
\midrule
Partecipanti & Matteo, Marco, Davide \\
\midrule
Redatto da & Generato da trascrizione automatica \\
\bottomrule
\end{tabular}
\end{center}

\vspace{0.8cm}

% ─────────────────────────────────────────────
\section{Argomenti Discussi}

\subsection{Architettura del sistema MirrorFish e analisi fondamentale}

Il team ha discusso la struttura del sistema MirrorFish, concepito come un modello linguistico dotato di ampie conoscenze generali ma privo di capacità di calcolo autonomo. Per ovviare a questa limitazione, il sistema necessita di ricevere dati pre-elaborati esterni, tra cui:

\begin{itemize}[leftmargin=1.5cm]
  \item \textbf{Analisi tecnica}: dati quantitativi, pre-section, e potenzialmente order flow.
  \item \textbf{Analisi fondamentale}: un insieme strutturato di informazioni finanziarie che Davide dovrà programmare e recuperare automaticamente.
\end{itemize}

È stato evidenziato che Davide non ha attualmente le competenze per definire \textit{come} svolgere l'analisi fondamentale. Il compito di Marco (e Matteo) sarà quello di redigere un documento guida che spieghi la procedura step-by-step, indicando quali dati raccogliere e come elaborarli per i diversi tipi di asset.

\subsection{Vincoli sui dati fondamentali}

È emerso un requisito fondamentale: i dati devono essere \textbf{reperibili storicamente in modo retroattivo}. Qualunque sia la fonte, il sistema deve essere in grado di ottenere i valori fondamentali riferiti a una data specifica del passato. Inoltre, la struttura della tabella di output deve essere definita \textbf{a priori e fissa}, per garantire la compatibilità con il contesto passato al modello.

\subsection{Organizzazione della documentazione del progetto}

È stata discussa la scelta dello strumento per gestire i documenti tecnici del progetto. Le opzioni valutate sono state:

\begin{itemize}[leftmargin=1.5cm]
  \item \textbf{Markdown}: facile da scrivere ma limitato nella navigazione tra più file e nella stampa in PDF strutturata.
  \item \textbf{Obsidian}: potente ma richiede configurazione e non si presta bene alla collaborazione.
  \item \textbf{LaTeX}: proposto come soluzione ottimale per creare documenti strutturati, compilabili in PDF, con capitoli, sezioni e sotto-sezioni facilmente navigabili nel file system.
  \item \textbf{Notion}: valutato come alternativa collaborativa, con sistema di link interni e accesso condiviso.
\end{itemize}

La proposta è quella di creare una \textbf{repository separata} dedicata esclusivamente alla documentazione, con file \LaTeX{} organizzati in cartelle per capitolo/sezione. I file possono essere scritti con l'ausilio di AI (es.\ ChatGPT o Claude) e compilati tramite Overleaf.

\vspace{0.5cm}

% ─────────────────────────────────────────────
\section{To Do}

\begin{enumerate}[leftmargin=1.5cm, label=\textbf{\arabic*.}]

  \item \textbf{Marco + Matteo — Documento guida sull'analisi fondamentale}\\
  Redigere un documento che spieghi a Davide come effettuare l'analisi fondamentale per i diversi tipi di asset finanziari. Il documento deve:
  \begin{itemize}
    \item Definire la procedura step-by-step;
    \item Specificare quali dati raccogliere e da quali fonti;
    \item Indicare le casistiche per tipo di asset (es.\ azioni, ETF, crypto, \dots);
    \item Essere comprensibile anche da chi non ha background economico (``deve capirlo anche Paolo'').
  \end{itemize}

  \item \textbf{Davide — Implementazione funzione fondamentale}\\
  Sulla base del documento guida, scrivere la funzione che:
  \begin{itemize}
    \item Raccoglie i dati fondamentali (scraping o API);
    \item Restituisce sempre la \textbf{stessa struttura tabellare} con campi definiti a priori;
    \item Supporta query storiche retroattive a una data specifica.
  \end{itemize}

  \item \textbf{Tutti — Setup repository documentazione}\\
  Creare una nuova repository GitHub dedicata ai documenti, con struttura a cartelle per capitolo. Valutare se usare \LaTeX{} (con Overleaf) o Notion, e invitare tutti i membri del team.

  \item \textbf{Tutti — Sessione di tutoraggio Overleaf/\LaTeX}\\
  Organizzare una sessione di onboarding per familiarizzare con la struttura dei file \LaTeX{} e il workflow di compilazione su Overleaf.

\end{enumerate}

\vspace{1cm}
\noindent\textcolor{darkgray}{\textit{Fine verbale.}}

\end{document}